% ALT TEXT: A schematic coordinate projection of one rhombic-dodecahedral cell contains a smaller centered circle representing its inscribed sphere. Four marked tangency points lie on the sphere and representative face planes, showing that the sphere touches the cell only at tangencies and does not coincide with its boundary. Labels record eta equals square root of 2, D equals pi over 3 square root of 2, Omega equals one minus D, and D plus Omega equals I.
\begin{figure}[t]
\centering
\begin{tikzpicture}[
  scale=1.05,
  every node/.style={font=\small},
  contact/.style={fill=red!70!black, draw=white, line width=.3pt}
]
\coordinate (O) at (0,0);
\coordinate (Ttop) at (0,.88);
\coordinate (TtopR) at (.42,.88);
\coordinate (Tright) at (.88,0);
\coordinate (TrightT) at (.88,.42);
\draw[very thick, blue!60!black]
  (0,2.25)--(1.15,1.55)--(2.25,.45)--(2.25,-.45)--
  (1.15,-1.55)--(0,-2.25)--(-1.15,-1.55)--(-2.25,-.45)--
  (-2.25,.45)--(-1.15,1.55)--cycle;
\draw[blue!35]
  (0,2.25)--(0,1.05)--(1.15,1.55)--(.88,.45)--(2.25,.45)
  (0,1.05)--(-.88,.45)--(-1.15,1.55)
  (0,-2.25)--(0,-1.05)--(1.15,-1.55)--(.88,-.45)--(2.25,-.45)
  (0,-1.05)--(-.88,-.45)--(-1.15,-1.55);
\fill[orange!18] (O) circle (.88);
\draw[very thick, orange!80!black] (O) circle (.88);
% Representative projected face planes tangent to the insphere.
\draw[red!60!black, thick] (-.42,.88)--(.42,.88);
\draw[red!60!black, thick] (-.42,-.88)--(.42,-.88);
\draw[red!60!black, thick] (.88,-.42)--(.88,.42);
\draw[red!60!black, thick] (-.88,-.42)--(-.88,.42);
\foreach \p in {(0,.88),(0,-.88),(.88,0),(-.88,0)}
  \node[contact,circle,inner sep=1.6pt] at \p {};
\draw[dashed, gray] (O)--(0,.88);
\draw[dashed, gray] (O)--(.88,0);
\draw pic[draw,angle radius=3mm] {right angle = O--Ttop--TtopR};
\draw pic[draw,angle radius=3mm] {right angle = O--Tright--TrightT};
\node[align=left, anchor=west] at (2.7,1.3)
  {\(\eta=\sqrt2\)\\[2pt]
   \(D=\dfrac{\pi}{3\sqrt2}\)\\[2pt]
   \(\Omega=1-D\)\\[2pt]
   \(\boxed{D+\Omega=I}\)};
\node[font=\scriptsize, text=orange!80!black] at (0,-.25)
  {centered insphere};
\node[font=\scriptsize, text=blue!60!black] at (0,2.55)
  {one RD cell};
\node[font=\scriptsize, text=red!60!black, align=center] at (3.55,-1.15)
  {representative orthogonal\\sphere--face tangencies};
\end{tikzpicture}
\caption{Coordinate presentation of the derived FCC/rhombic-dodecahedral
cell, its centered inscribed sphere, and representative orthogonal tangencies.
The sphere and cell boundary are distinct except at the tangency points.  The
drawing presents derived relational geometry; it is not a primitive external
lattice or background space.}
\label{fig:fcc-rd-cell}
\end{figure}
